% ─────────────────────────────────────────────────────────────────────────────
%  preamble.tex  –  Paquetes y configuración global
% ─────────────────────────────────────────────────────────────────────────────

% ── Codificación y lengua ────────────────────────────────────────────────────
\usepackage[utf8]{inputenc}
\usepackage[T1]{fontenc}
\usepackage[spanish,es-noshorthands]{babel}

% ── Tipografía ───────────────────────────────────────────────────────────────
\usepackage{lmodern}
\usepackage{microtype}

% ── Matemáticas ──────────────────────────────────────────────────────────────
\usepackage{amsmath,amssymb,amsthm}
\usepackage{mathtools}

% ── Geometría de página ──────────────────────────────────────────────────────
\usepackage[
  a4paper,
  top=2.5cm, bottom=2.5cm,
  left=3cm,  right=2.5cm,
  headheight=15pt
]{geometry}

% ── Gráficos y diagramas ─────────────────────────────────────────────────────
\usepackage{graphicx}
\usepackage{tikz}
\usetikzlibrary{shapes.geometric, arrows.meta, positioning, fit, calc, backgrounds}
\usepackage{pgfplots}
\pgfplotsset{compat=1.18}

% ── Código fuente ────────────────────────────────────────────────────────────
\usepackage{listings}
\usepackage{xcolor}

\definecolor{codegreen}{rgb}{0.13,0.55,0.13}
\definecolor{codegray}{rgb}{0.5,0.5,0.5}
\definecolor{codepurple}{rgb}{0.58,0,0.82}
\definecolor{codeblue}{rgb}{0.0,0.2,0.7}
\definecolor{backcolour}{rgb}{0.97,0.97,0.97}
\definecolor{keywordcolor}{rgb}{0.0,0.2,0.7}

\lstdefinestyle{java}{
  backgroundcolor=\color{backcolour},
  commentstyle=\color{codegreen},
  keywordstyle=\color{keywordcolor}\bfseries,
  numberstyle=\tiny\color{codegray},
  stringstyle=\color{codepurple},
  basicstyle=\ttfamily\footnotesize,
  breakatwhitespace=false,
  breaklines=true,
  captionpos=b,
  keepspaces=true,
  numbers=left,
  numbersep=5pt,
  showspaces=false,
  showstringspaces=false,
  showtabs=false,
  tabsize=4,
  language=Java,
  morekeywords={@pre,@post,@invariant,@variant,result},
  frame=single,
  rulecolor=\color{gray!40},
}

\lstdefinestyle{asm}{
  backgroundcolor=\color{backcolour},
  commentstyle=\color{codegreen}\itshape,
  keywordstyle=\color{codeblue}\bfseries,
  basicstyle=\ttfamily\footnotesize,
  breaklines=true,
  captionpos=b,
  numbers=left,
  numberstyle=\tiny\color{codegray},
  numbersep=5pt,
  frame=single,
  rulecolor=\color{gray!40},
  language={[x86masm]Assembler},
}

\lstdefinestyle{smt}{
  backgroundcolor=\color{backcolour},
  commentstyle=\color{codegreen}\itshape,
  keywordstyle=\color{codepurple}\bfseries,
  basicstyle=\ttfamily\footnotesize,
  breaklines=true,
  captionpos=b,
  numbers=left,
  numberstyle=\tiny\color{codegray},
  numbersep=5pt,
  frame=single,
  rulecolor=\color{gray!40},
  morekeywords={declare-fun,assert,check-sat,get-model,push,pop,exit,set-logic,not,and,or,=>,>=,<=,+,-,*,/},
}

\lstdefinestyle{cpp}{
  backgroundcolor=\color{backcolour},
  commentstyle=\color{codegreen}\itshape,
  keywordstyle=\color{keywordcolor}\bfseries,
  basicstyle=\ttfamily\footnotesize,
  breaklines=true,
  captionpos=b,
  numbers=left,
  numberstyle=\tiny\color{codegray},
  numbersep=5pt,
  frame=single,
  rulecolor=\color{gray!40},
  language=C++,
}

\lstset{style=java}   % estilo por defecto

% ── Tablas ───────────────────────────────────────────────────────────────────
\usepackage{booktabs}
\usepackage{tabularx}
\usepackage{multirow}
\usepackage{array}

% ── Hiperenlaces ─────────────────────────────────────────────────────────────
\usepackage[
  colorlinks=true,
  linkcolor=black,
  citecolor=codeblue,
  urlcolor=codeblue,
  pdftitle={Compilador JavaSubset → x86-64 con Verificación Formal},
  pdfauthor={nickhernd},
]{hyperref}

% ── Encabezados y pies de página ─────────────────────────────────────────────
\usepackage{fancyhdr}
\pagestyle{fancy}
\fancyhf{}
\fancyhead[LE,RO]{\thepage}
\fancyhead[RE]{\small\nouppercase{\leftmark}}
\fancyhead[LO]{\small\nouppercase{\rightmark}}
\renewcommand{\headrulewidth}{0.4pt}

% ── Entornos matemáticos ─────────────────────────────────────────────────────
\theoremstyle{definition}
\newtheorem{definicion}{Definición}[chapter]
\newtheorem{ejemplo}{Ejemplo}[chapter]
\newtheorem{regla}{Regla}[chapter]

\theoremstyle{plain}
\newtheorem{teorema}{Teorema}[chapter]
\newtheorem{lema}{Lema}[chapter]

\theoremstyle{remark}
\newtheorem{nota}{Nota}[chapter]
\newtheorem{observacion}{Observación}[chapter]

% ── Pseudocódigo (sin paquete algorithm2e) ───────────────────────────────────
\lstdefinestyle{pseudo}{
  backgroundcolor=\color{backcolour},
  basicstyle=\ttfamily\footnotesize,
  breaklines=true,
  captionpos=b,
  frame=single,
  rulecolor=\color{gray!40},
  numbers=left,
  numberstyle=\tiny\color{codegray},
  numbersep=5pt,
}

% ── Miscelánea ───────────────────────────────────────────────────────────────
\usepackage{enumitem}
\usepackage{caption}
\usepackage{subcaption}
\usepackage{float}
\usepackage{csquotes}
% backnaur no disponible; gramáticas se presentan en lstlisting

% ── Macros del proyecto ──────────────────────────────────────────────────────
\newcommand{\jss}{\texttt{JavaSubset}}
\newcommand{\IR}{\textsc{ir}}
\newcommand{\WP}[2]{\mathrm{WP}(#1,\,#2)}
\newcommand{\SP}[2]{\mathrm{SP}(#1,\,#2)}
\newcommand{\hoare}[3]{\{#1\}\;#2\;\{#3\}}
\newcommand{\subst}[3]{#1[#3/#2]}           % P[e/x]
\newcommand{\true}{\mathtt{true}}
\newcommand{\false}{\mathtt{false}}
\newcommand{\reg}[1]{\texttt{\%#1}}
